% !TEX program = xelatex
\documentclass[11pt,a4paper]{ctexart}
\usepackage[margin=2.2cm]{geometry}
\usepackage{amsmath,amssymb}
\usepackage{xcolor}
\usepackage{enumitem}
\usepackage{booktabs}
\usepackage{longtable}
\usepackage{hyperref}
\hypersetup{colorlinks=true,linkcolor=blue!50!black}

\setlength{\parskip}{3pt}
\newcounter{finding}
\newenvironment{finding}[1]{%
  \refstepcounter{finding}%
  \subsection*{条目 \thefinding\quad #1}%
  \addcontentsline{toc}{subsection}{条目 \thefinding\quad #1}%
}{\par\medskip\noindent\rule{\linewidth}{0.3pt}}
\newcommand{\loc}[1]{\par\noindent\textbf{【原文位置】}#1}
\newcommand{\orig}[1]{\par\noindent\textbf{【原文内容】}\begin{quote}\small\itshape #1\end{quote}}
\newcommand{\ptype}[1]{\par\noindent\textbf{【问题类型】}#1}
\newcommand{\basis}[1]{\par\noindent\textbf{【具体依据及对理解或论证的影响】}#1}
\newcommand{\conf}[1]{\par\noindent\textbf{【判断把握】}#1}
\newcommand{\normal}[1]{\par\noindent\textbf{【是否可能属于正常学术表达及原因】}#1}
\newcommand{\advice}[1]{\par\noindent\textbf{【建议】}#1}
\newcommand{\revise}[1]{\par\noindent\textbf{【建议修改文本】}\begin{quote}\small #1\end{quote}}
\newcommand{\authfill}[1]{\textcolor{red!70!black}{[待作者补充：#1]}}
\newcommand{\verify}{\textcolor{orange!80!black}{\textbf{待核实}}}

\title{M4Fed 稿件 AI 生成文本特征与语言质量审查报告}
\author{审查对象：\texttt{latex/main.tex} 与 \texttt{latex/supplementary.tex}}
\date{2026年11月6日}

\begin{document}
\maketitle

\begin{abstract}
本报告对稿件 \emph{A Multi-Source Federated Learning Framework with Medical Foundation Models for Multi-Modal Medical Data on Brain Age Estimation}（正文 \texttt{main.tex}，补充材料 \texttt{supplementary.tex}）做写作特征与语言质量审查。审查只分析文本的写作特征、论证与证据之间的对应关系，以及文内一致性，用于改进稿件，不判断文本是否由 AI 生成，也不给出任何“AI 生成概率”。文中“第 N 行”均指 \texttt{main.tex} 源文件行号（补充材料记为“S 第 N 行”），与 PDF 中的行号不一定一致。所有数值核对都只用稿件自身给出的数字和图；引用文献只核对了 \texttt{references.bib} 中的题名，没有检索全文，凡涉及“引用是否支持主张”的判断均标为“\verify”。
\end{abstract}

\tableofcontents
\newpage

%=====================================================================
\section{整体观察}

\subsection{材料与主要写作特点}
稿件提出 M4Fed 框架（TabPFN 补全 + MedMD 文本条件 3D 潜扩散 + SwinMTL 多任务预测，客户端使用 QLoRA），并在 OpenMind T1w 数据上做脑龄估计实验。全文结构完整，实验部分信息量大，数值较多。总体写作特点如下：

\begin{enumerate}[leftmargin=*]
\item \textbf{“三元”框架反复复述。}“三类资源（私有数据、公开数据、基础模型）”与“三类约束（隐私、通信、边缘算力）”在摘要、引言（第102、104、106、111行）、相关工作（第159行）、问题定义（第188行）、方法概述（第200行）、讨论（第684、694行）和结论（第703行）中至少出现八次，措辞高度相似。第694行“Practical implications”段与第703行结论基本是引言的换词复述，没有增加新信息。这是本稿最集中的写作问题（问题类型2、6）。
\item \textbf{结论表述的范围大于证据。}摘要、引言、结论中的“surpasses state-of-the-art”“reduce communication payload by 42.6\%”“simultaneously improve accuracy, reduce communication and lower edge compute”在正文中都有附加条件：只有服务器端优于基线；42.6\% 只在非默认的 $\rho_f=0.01$ 下成立，默认 $\rho_f=0.1$ 时上传量是 FedAvg 的 5.55 倍；边缘算力没有实测，只统计了可训练参数量。条件在第5.8节写得较清楚，但没有回传到摘要、表1和结论（问题类型3、9）。
\item \textbf{机制解释较多，但均为未经验证的推测。}结果讨论中，“We see two plausible reasons”“A possible / plausible explanation is that”等句式反复出现（第455、517、574、593、604、643、645、686行）。多数段落采用“结果 → 两条可能原因 → 结论性总结”的固定展开模式，有些段落在列出推测后又以确定语气作结（如第604、686行），把推测写成了结论（问题类型2、7）。
\item \textbf{文内一致性问题集中在“默认配置结果”和“评估单位”。}同一默认配置下的 M4Fed 服务器结果在三张表中分别是 MAE 8.87、9.376、9.22；主比较把单个服务器节点与基线的客户端均值对比；数据划分单位（扫描或受试者）没有交代（问题类型8、9）。
\item \textbf{做得较好的部分。}第5.8节通信量推算的数值全部可复算且自洽，并主动指出默认配置下 M4Fed 的上传量反而更大；第5.9节集中式实验的统计检验写明了检验方法、样本量和 $p$ 值，对小样本格子也做了提醒；局限性部分（第692行）主动承认了模拟环境和隐私方面的不足。这些段落可以作为全文修改的参照。
\end{enumerate}

\subsection{数值一致性核对（已通过的项）}
以下数值经复算与原文一致，不需要修改：队列 $3{,}622+4{,}571-875=7{,}318$；$875/3{,}622\approx24.2\%$，与未筛选队列中的健康比例一致；表3 合并行各百分比；第448行 $13.7\%$、$1.8\times$；第455行 $5.4$--$5.6\%$；第480行 $74\%$、$21\%$；第517行 $2.0\%$；第574行注册实验的所有百分比；第602行 $6.2\%$；第645行 $-20.8\%$、$-32.4\%$；第651--653行通信与参数量的全部推算（$9{,}247{,}036$、$1{,}669{,}770$、$81.9\%$、$92.6\%$、$205.3$~MB、$21.2$~MB、$5.55\times$、$42.6\%$、$\rho_f\approx0.018$、$19.5\times$）；图 centralized 中女性 307 + 男性 256 + 缺失 8 $=571$；图 test\_scatter\_registered 中三节点测试样本 $188+214+169=571$，客户端平均 MAE $(5.26+4.69)/2\approx4.97$，与表 registration 一致。跨站点评估中，由 RMSE 与 $R^2$ 反推的年龄标准差（$13.99/\sqrt{1.114}\approx13.3$；$14.36/\sqrt{0.822}\approx15.8$）也与第602行给出的 $13.3$ 与 $15.9$ 年一致。

\subsection{问题是否集中出现}
问题集中在四个位置：(1) 摘要、引言末段、表1、讨论末段和结论中的总括性论断；(2) 第5.2--5.7节结果讨论里的“推测性解释→确定性结论”过渡；(3) 隐私相关表述；(4) 方法与实验设置中缺失的复现信息。方法细节段落（第4.2--4.4节）以技术描述为主，语言问题较少，主要问题是个别超长句和未定义术语。

%=====================================================================
\section{逐条定位}

%---------------------------------------------------------------------
\begin{finding}{摘要中的样本量表述与实际实验队列不符}
\loc{摘要，第84行；对照表3（第349--363行）、第5.4、5.6、5.9节。}
\orig{To train and validate M4Fed, we performed brain-age estimation on 7318 subjects of public datasets.}
\ptype{9 内容一致性；3 结论范围放大。}
\basis{7{,}318 是两个队列去重合并后的人数（表3“Deduplicated combine”行）。实际实验分别在未筛选队列（3{,}622 人）和健康对照队列（4{,}571 人）上进行，没有任何一组实验同时使用 7{,}318 人。此外，正文没有明确说明主比较（表5）、消融（表6）和敏感性分析（表8）用的是哪个队列，只能从表7的说明中间接推测为未筛选队列。读者会误以为主要结果基于 7{,}318 人。}
\conf{高}
\normal{用合并人数概括数据规模是常见写法，但前提是正文说明各实验分别使用的子集。本稿正文没有交代主实验所用队列，所以这里属于实质性的信息缺口。}
\advice{改写；并在第5.1或5.2节明确每组实验所用的队列。}
\revise{We evaluate M4Fed on brain-age estimation using two cohorts drawn from the public OpenMind corpus: an unselected cohort (3{,}622 subjects) and a healthy-controls cohort (4{,}571 subjects), 7{,}318 unique subjects in total. \authfill{说明主比较、消融与敏感性实验分别使用哪个队列}}
\end{finding}

%---------------------------------------------------------------------
\begin{finding}{摘要、引言、结论中的性能与通信结论缺少成立条件}
\loc{摘要第84行；引言第141行；结论第703行；对照表5（第429--446行）、表8（第519--540行）、第5.8节（第651--653行）、第688行。}
\orig{M4Fed surpasses state-of-the-art approaches while reducing communication payload by 42.6\% and edge-side parameter burden by over 80\%.\\
(第141行) The experimental results show that M4Fed outperforms the other baselines, while also reducing communication payload by 42.6\% ...\\
(第703行) M4Fed can simultaneously improve predictive accuracy, reduce communication overhead, and lower edge-side compute.}
\ptype{3 结论被放大；9 内容一致性。}
\basis{(a) 精度：只有 M4Fed 服务器在全部指标上最优；M4Fed 客户端的 MAE（10.97）、RMSE（17.19）不优于基线，$R^2$（$-0.184$）是全表最低（表5，第448行亦承认）。(b) 通信：42.6\% 只在 $\rho_f=0.01$、单个客户端（$c_1$）、单轮上传时成立；默认 $\rho_f=0.1$（表4，也是表5所用配置）下上传量是 FedAvg 的 $5.55$ 倍（第653、688行）。因此摘要中的精度和通信两项结论来自不同配置，而 $\rho_f=0.01$ 下的精度只在表8中有单次结果。(c) 统计范围：比较只计入客户端上传，没有计入服务器→客户端下发的合成数据、一次性 $|\theta_{\mathrm{GM}}|$ 广播（30.1~MB），也没有计入 FedAvg 的下行量。(d) “parameter burden over 80\%”指域模型可训练参数减少 81.9\%，并非计算量或显存；客户端还需额外运行扩散生成（第692行局限 iv），结论中的“lower edge-side compute”缺少实测支撑。(e) 基线包含 FedAvg（2017）和 FedProx（2020），统称“state-of-the-art approaches”不准确。}
\conf{高}
\normal{摘要压缩信息属正常做法，但省略的恰好是决定结论正负的条件（默认配置下通信反而增加），已超出正常压缩范围。}
\advice{改写；摘要、引言、结论三处同步修改。}
\revise{(摘要) The M4Fed server model attains the lowest MAE among the compared federated methods (8.87 vs.\ 10.28 years for the best baseline), whereas the M4Fed client models do not improve point-error metrics over the baselines. QLoRA reduces the client-side trainable parameters of the domain model by 81.9\%. For the 9.25M-parameter domain model used here, the per-round client upload is 42.6\% smaller than that of FedAvg at a synthesis ratio of $\rho_f=0.01$, but 5.55 times larger at the default $\rho_f=0.1$.\\
(结论) ... M4Fed improves server-side accuracy over the compared baselines and reduces client-side trainable parameters; its communication advantage holds only below a synthesis ratio of about 0.018 for the compact model studied, and is expected to grow with model size.}
\end{finding}

%---------------------------------------------------------------------
\begin{finding}{方法名称的展开与缩写不对应；摘要中的对应关系错位}
\loc{题目；摘要第84行；第111行。}
\orig{we propose a \textbf{M}ulti-source \textbf{Fed}erated Learning with \textbf{M}ulti-\textbf{M}odal \textbf{M}edical \textbf{F}oundation-model (M4Fed) method.\\
Consolidating these resources however faces obstacles, including 1) privacy ..., 2) ... communication costs, and 3) limited computing resources ... M4Fed 1) connects ..., 2) offloads edge-side computation ..., 3) employs synthetic-data-driven federated aggregation to keep communication cost manageable.}
\ptype{9 一致性；5 表达不自然。}
\basis{加粗字母依次为 M、Fed、M、M、M、F，与“M4Fed”（四个 M）对不上；第111行的全称没有加粗，读者无法得知命名依据。摘要中障碍 2)/3)（通信/算力）与对策 2)/3)（算力/通信）顺序颠倒，读者需要自行对应。“Consolidating these resources however faces obstacles, including 1) privacy regulations restrict ...”中“including”后接完整从句，语法不当；“large-scale parameters exchanging”搭配不自然。}
\conf{中}
\normal{缩写不必与全称逐字母对应，但本稿刻意加粗字母，读者会按加粗字母去对应，因此需要统一。}
\advice{改写。}
\revise{Consolidating these resources faces three obstacles: (1) privacy regulations restrict the centralization of patient data; (2) exchanging large parameter sets incurs high communication costs; and (3) hospital nodes lack the compute for full-scale local training. M4Fed addresses (1) through federated learning across isolated repositories, (2) through synthetic-data-driven aggregation, and (3) through parameter-efficient fine-tuning with cloud-edge collaboration. \authfill{给出与“M4Fed”逐字对应的全称，或删除加粗字母}}
\end{finding}

%---------------------------------------------------------------------
\begin{finding}{研究动机依赖“基础模型规模”，但实验模型规模较小}
\loc{第104行；第156行；第313行；对照第651、653行。}
\orig{(第104行) not for the large, volumetric foundation models that are now increasingly influencing medical image analysis ...\\
(第156行) ... once the aggregated model is a volumetric foundation network with tens or hundreds of millions of parameters [ren2025fedfm], as is now standard in medical image analysis ...\\
(第313行) ... whose message payloads become prohibitively large under full-model federated protocols once the shared model reaches foundation scale.}
\ptype{3 结论放大；9 一致性；7 证据边界。}
\basis{引言和方法以“foundation-scale 参数导致通信不可承受”为核心动机，但实验中的域模型为 9.25M 参数，生成模型的可训练部分为 7.54M（第651行）。在这一规模下，默认配置的通信量高于 FedAvg（第653行）。1B 参数情形只是推算，没有实验。“as is now standard in medical image analysis”是概括性论断；所引 ren2025fedfm 的题名为 \emph{Advances and Open Challenges in Federated Foundation Models}，是否支持“在医学影像中已成为标准”\verify。“prohibitively large”没有量化。}
\conf{高}
\normal{用领域趋势引出动机是常见写法。问题在于动机与实验规模之间的差距直到第5.8节才说明，前文的表述会让读者高估实验对动机的验证程度。}
\advice{改写；在引言或第5.1节提前说明实验模型规模，并区分实测与推算。}
\revise{(第156行后半句) ... becomes a bottleneck as model size grows towards that of foundation models \authfill{引用支持“医学影像基础模型参数规模”的具体文献}. \\
(第5.1节新增) In all experiments the domain model has 9.25M parameters; communication costs for larger, foundation-scale models are extrapolated analytically in Section~5.8 rather than measured.}
\end{finding}

%---------------------------------------------------------------------
\begin{finding}{引言第三段：修饰性表达堆叠、单复数不一致}
\loc{第106行。}
\orig{In recent years, a series of methods proposed by researchers have gradually been converging into a technical line that may be collectively termed multi-source federated learning, progressively integrating training resources scattered across disparate sources. ... This alleviates the predicament in which edge nodes, constrained by limited computational resources, struggle to fully fine-tune massive foundation models, and it thereby forges an effective link between foundation models and private data. Finally, some of the latest studies [zhang2025gptfl] employ a large language model ... Before entering the federated loop, it trains the downstream model ... thereby further mobilizing the knowledge embedded within foundation models in the form of synthetic data, going beyond the private data alone.}
\ptype{5 表达不自然；1 空泛；2 机械重复。}
\basis{“a series of methods proposed by researchers have gradually been converging into a technical line”“alleviates the predicament”“forges an effective link”“further mobilizing the knowledge embedded within”等属于比喻性、评价性表达，没有提供可检验的信息。“thereby”在一段内出现三次，“First/Second/Third/Finally”结构与上一段（第104行）重复。“some of the latest studies”只引用一篇文献，后句又用单数“it”。本段是对相关工作的预告，与第2.1节第二段内容重复。}
\conf{中}
\normal{综述式铺垫本身属正常学术写法；但比喻性动词和“thereby”连用使论证显得冗长，可以精简而不损失信息。}
\advice{精简。}
\revise{Recent work has progressively enlarged the set of training resources available to federated learning, a line we refer to as multi-source federated learning. Classical federated learning [FedAvg, FedProx] learns from decentralized private data by aggregating model weights instead of raw records. FedDF [lin2020feddf] adds unlabeled public data through ensemble distillation. FedPEFT and OpenFedLLM [sun2022fedpeft, ye2024openfedllm] fine-tune only low-rank adapters of a frozen foundation model, which makes such models trainable on resource-limited clients. GPT-FL [zhang2025gptfl] uses a large language model and a generative model to synthesize labeled data for server-side warm-up before the federated loop.}
\end{finding}

%---------------------------------------------------------------------
\begin{finding}{相关工作正文与表1对现有方法的定位相互矛盾}
\loc{第159行；表1（第113--138行）；第111行。}
\orig{Each of these methods, however, enlarges the resource pool along a single axis only. None combines private client data, public repositories, and pretrained foundation models within one framework ...}
\ptype{9 内容一致性；8 文献综述。}
\basis{表1中 GPT-FL 在 Private / Public / Foundation 三列都标为 \checkmark，FedPEFT、OpenFedLLM 也标了两类资源，与“single axis only”“None combines ...”直接矛盾。第111行的表述（“at most one or two of these constraints and draws on only a subset”）与表1一致，与第159行不一致。另外，GPT-FL 在表1中标为使用 public data，依据似乎是本文复现中“GPT-FL's server reuses the generator trained by M4Fed's”（第367行）。原论文是否使用公开数据集\verify。若是复现设置带来的，应在表注中说明。FedPDM 在正文中讨论，但没有出现在表1中。}
\conf{高}
\normal{否。属于可直接核对的事实矛盾。}
\advice{改写，并核实表1中 GPT-FL 一行的标注依据。}
\revise{Each of these methods enlarges the resource pool in some directions, but none of them mitigates the privacy, communication, and edge-compute constraints at the same time while drawing on private data, public data, and foundation models (Table~1).}
\end{finding}

%---------------------------------------------------------------------
\begin{finding}{表1对 M4Fed 的“Communication”“Edge compute”标注无条件对勾}
\loc{表1第133行；对照第653、688、692行。}
\orig{M4Fed (ours) \& \checkmark \& \checkmark \& \checkmark \& \checkmark \& \checkmark \& \checkmark}
\ptype{3 结论放大；9 一致性。}
\basis{第653行表明，在本文模型规模和默认 $\rho_f=0.1$ 下，M4Fed 的单轮上传量是 FedAvg 的 5.55 倍；第692行（局限 iv）承认生成模型会增加节点的训练成本和显存。表1对其他方法使用了 $\triangle$ 加脚注的条件标注，对本方法却不加条件，评价标准前后不一致。}
\conf{高}
\normal{定位表中对本方法打满对勾很常见，但本稿自己的结果表明这两项是有条件成立的。}
\advice{改写表格标注或增加脚注。}
\revise{Communication 与 Edge compute 两列改为 $\triangle$，并加脚注：“Per-round payload is decoupled from model size and is smaller than FedAvg only when $\rho_f n_k\beta < 4|\theta_{\mathrm{DM}}|$ (Section~5.8); client-side trainable parameters are reduced by QLoRA, but clients additionally run diffusion sampling.”}
\end{finding}

%---------------------------------------------------------------------
\begin{finding}{“knowledge diffuses across institutions”与协议实际信息流不一致；域模型不参与联邦聚合未说明}
\loc{第111行 (i)；第204行；算法1第239--247行；第188行约束(c)。}
\orig{(i) ... allowing knowledge to diffuse across institutions without exposing raw data.\\
(第188行) (c) the compute available to fine-tune $\theta$ at each client $c_k$ is insufficient to backpropagate through the full parameter set of $f_\theta$ at the frequency required by the federated protocol.}
\ptype{4 逻辑跳跃；9 一致性。}
\basis{根据算法1，SwinMTL（域模型）在 M4Fed 中从未被聚合或交换，各节点只在第3阶段用本地私有数据和本地合成数据训练。跨节点信息只能经由 MedMD 的合成数据间接传递，而客户端的 MedMD 又只在服务器下发的合成数据上微调。正文没有把这一点明确告诉读者，使“knowledge diffuses across institutions”的含义模糊；约束 (c) 中“at the frequency required by the federated protocol”对 M4Fed 的域模型并不适用。这也有助于解释为何客户端结果不优于基线（表5）。}
\conf{中}
\normal{属于描述精度问题，但它影响读者对客户端结果的理解。}
\advice{改写，在第4.1节明确信息流。}
\revise{In M4Fed only the generative model is trained collaboratively; the domain model SwinMTL is never aggregated and is fine-tuned locally at each node on its own private and synthetic data. Cross-institution information therefore reaches a node's domain model only through the synthetic volumes produced by its generative model.}
\end{finding}

%---------------------------------------------------------------------
\begin{finding}{贡献表述中“comprehensive”缺少支撑}
\loc{第146行。}
\orig{We conduct a comprehensive empirical evaluation against classical and state-of-the-art federated learning methods on OpenMind, complemented by ablation, sensitivity, and cross-site generalization studies.}
\ptype{3 重要性放大。}
\basis{实验为单次运行（文中没有报告随机种子或重复次数），M4Fed 只有 2 个客户端，没有显著性检验，跨站点评估只评估了 M4Fed 服务器，没有评估基线。“comprehensive”与这些限制不相称。}
\conf{中}
\normal{“comprehensive”在贡献列表中较常见，单独看属轻微问题；放在本稿整体结论偏强的背景下，建议去掉。}
\advice{精简。}
\revise{We evaluate M4Fed against four federated baselines on OpenMind, together with component ablations, sensitivity analyses, and a cross-site evaluation of the server model.}
\end{finding}

%---------------------------------------------------------------------
\begin{finding}{相关工作第一段的超长句}
\loc{第156行。}
\orig{Both protocols transmit the full parameter set of the model being trained at every communication round, an assumption that is largely benign for compact convolutional classifiers but becomes a first-order bottleneck once the aggregated model is a volumetric foundation network with tens or hundreds of millions of parameters [ren2025fedfm], as is now standard in medical image analysis; model compression, gradient sparsification, evolutionary knowledge-transfer, and asynchronous update strategies have all been proposed as partial remedies for this bottleneck [...].}
\ptype{5 超长句；4 概念堆砌。}
\basis{该句约 80 词，包含同位语、两个从句和分号连接的独立句，涵盖问题、规模、领域现状、四类补救方法四层信息，读者难以把握重点。“a first-order bottleneck”为修辞性表达。}
\conf{中}
\normal{长句是学术英语的常见风格，但本句已影响理解，拆分不会损失信息。}
\advice{改写（拆句）。}
\revise{Both protocols transmit the full parameter set at every round. This is affordable for compact classifiers but becomes a bottleneck for large volumetric models. Model compression, gradient sparsification, evolutionary knowledge transfer, and asynchronous updates have been proposed to reduce this cost [wang2025fedect, sabah2025communication]. Because all of these act on how the parameters are transmitted, the payload remains coupled to model size.}
\end{finding}

%---------------------------------------------------------------------
\begin{finding}{对 Swin UNETR 预训练与 TabPFN 的描述混合了文献结论与本文用法}
\loc{第163行。}
\orig{Tang et al. pretrain a hierarchical Swin Transformer encoder on about 5,000 unlabeled 3D computed tomography volumes, yielding representations that transfer strongly to downstream segmentation and, as we exploit in this work, regression tasks ...\\
Its most recent revision performs accurate regression and classification on small, heterogeneous clinical tables with no per-dataset hyperparameter tuning ...}
\ptype{7 事实/证据边界；8 引用对应。}
\basis{“transfer strongly to ... regression tasks”把本文的用法（作者自己的选择）与文献的发现（分割）并列成“已证实”的结论；本文消融（表6）中去掉预训练后 MAE 仅变化约 0.6 年，且是单次运行，不足以说明“strongly”。CT 预训练迁移到 T1w MRI 存在模态差异，本文没有讨论。TabPFN 文献是否专门针对“clinical tables”\verify（bib 题名不含 clinical）。}
\conf{中}
\normal{“as we exploit in this work”提示了这是本文用法，但与“transfer strongly”放在同一谓语下，容易被读成文献结论。}
\advice{改写，分开表述文献证据与本文用法。}
\revise{Tang et al. [tang2022swinunetr] pretrain a hierarchical Swin Transformer encoder on about 5{,}000 unlabeled 3D CT volumes and report gains on downstream segmentation. In this work we use these weights to initialize a regression model for T1-weighted MRI, a setting that differs from the pretraining data in both modality and task.}
\end{finding}

%---------------------------------------------------------------------
\begin{finding}{生成模型相关工作为罗列式综述，缺少比较与综合；缺少脑龄估计文献综述}
\loc{第167行；第2节整体。}
\orig{Diffusion models [...] is a paradigm for high-quality conditional image synthesis. Latent diffusion [...] relocates the diffusion process ..., reducing the compute and memory footprint of volumetric synthesis to a level compatible with hospital-grade hardware. More recently Kazerouni et al. ... Trinh et al. show ... Then Mirowski et al. instantiate this direction concretely ...}
\ptype{8 文献综述与引用；5 语法。}
\basis{本段按“More recently … Then …”逐篇罗列，没有比较已有 3D 脑 MRI 扩散生成方法与 MedMD 的差异（例如条件方式、分辨率、是否用于联邦场景）。“to a level compatible with hospital-grade hardware”没有依据。“Diffusion models ... is”主谓不一致。第2节没有脑龄估计方法（包括联邦脑龄估计，如第141行所引 denissen2026federated）的综述，而脑龄是本文的唯一应用任务，读者无法判断 MAE 8.87 年在该领域处于什么水平。}
\conf{中}
\normal{相关工作逐篇介绍是常见写法；但缺少与本文方法的对比，以及缺少应用任务的文献，属于实质性不足。}
\advice{改写并补充文献（由作者提供，本报告不新增文献）。}
\revise{Diffusion models [ho2020ddpm, song2021ddim, nichol2021improved] are widely used for conditional image synthesis, and latent diffusion [rombach2022ldm] reduces their cost by operating in the latent space of a pretrained autoencoder. \authfill{与已有 3D 脑 MRI 生成方法逐项比较：条件信号、分辨率、是否用于数据交换}. Unlike these works, MedMD conditions a 3D latent diffusion U-Net on BiomedCLIP embeddings of imputed metadata and is used as the communication medium of a federated protocol.\\
\authfill{新增“Brain age estimation”小节：代表性方法、常用数据集上的 MAE 范围、联邦脑龄估计工作}}
\end{finding}

%---------------------------------------------------------------------
\begin{finding}{PEFT 相关工作末尾的套话式过渡}
\loc{第171行。}
\orig{These efforts collectively affirm the effectiveness of parameter-efficient tuning for edge nodes within federated settings. Therefore, M4Fed applies it to a 3D vision backbone and to a diffusion U-Net ...}
\ptype{1 空泛与套话；2 机械过渡。}
\basis{“collectively affirm the effectiveness”是概括性评价，没有说明在何种条件、何种指标下有效；“Therefore”把文献结论直接当作本文设计的理由，逻辑过渡较生硬。前文所引工作多为 2D 或语言模型，本文用于 3D 骨干与扩散 U-Net，适用性的差异没有说明。}
\conf{中}
\normal{综述段落以总结句收尾属常见写法；建议改为说明差异，而不是笼统肯定。}
\advice{改写。}
\revise{These works apply parameter-efficient tuning to language models and 2D image classifiers. M4Fed applies QLoRA to a 3D vision backbone and a 3D diffusion U-Net at the clients, while the server performs full-parameter fine-tuning.}
\end{finding}

%---------------------------------------------------------------------
\begin{finding}{问题定义中的 $f_\theta(x,m)$ 与方法中“仅由 MRI 预测”不一致}
\loc{第185行式(2)；第288行；题目中的“Multi-Modal”。}
\orig{(式2) $\mathcal{L}_{\mathrm{MT}}(f_\theta(x,m), y)$\\
(第288行) The domain model SwinMTL predicts age and a set of auxiliary clinical fields from MRI volumes.}
\ptype{9 内容一致性；4 概念含糊。}
\basis{式(2)表明域模型以影像 $x$ 和元数据 $m$ 为输入；第4.4节描述的 SwinMTL 只以影像为输入，元数据是辅助预测目标。若推断时不使用元数据，题目和全文中的“multi-modal”只体现在生成模型的条件输入上，应明确说明；若推断时使用元数据，第4.4节缺少融合方式的描述。}
\conf{高}
\normal{否，属于形式化定义与实现之间的不一致。}
\advice{核实后改写。}
\revise{\authfill{确认 SwinMTL 推断时是否使用元数据}. 若仅用影像：将式(2)改为 $\mathcal{L}_{\mathrm{MT}}(f_\theta(x), y)$，并补充 “Metadata enter M4Fed in two ways: as text conditioning of the generative model and as auxiliary prediction targets of the domain model; the domain model itself takes only the MRI volume as input.”}
\end{finding}

%---------------------------------------------------------------------
\begin{finding}{术语“granularity”与“privacy-friendly”未定义}
\loc{第188行；第200行；第106行。}
\orig{(第188行) M4Fed satisfies ... (b) by the adjustable synthesis quantity and granularity instead of scaling with $|\theta|$ ...\\
(第200行) ... synthetic-data aggregation whose volume and granularity are user-configurable.\\
(第106、200行) ... in a privacy-friendly manner ...}
\ptype{4 含糊模糊、概念堆砌。}
\basis{全文只定义了合成比例 $\rho_f,\rho_l$ 与噪声强度 $\sigma$，没有任何参数对应“granularity”。“privacy-friendly”不是有明确含义的隐私术语，本文也没有进行隐私评估，该词容易被读成对隐私保护程度的承诺。}
\conf{中}
\normal{“privacy-friendly”偶见于文献，但在没有隐私评估的前提下宜采用中性描述。}
\advice{改写。}
\revise{(b) by exchanging a number of synthetic volumes set by the synthesis ratio $\rho_f$, which does not scale with $|\theta|$ ... \authfill{若“granularity”另有所指，请给出定义与对应参数}；“in a privacy-friendly manner” 改为 “without transferring raw images between nodes”.}
\end{finding}

%---------------------------------------------------------------------
\begin{finding}{隐私相关表述的证据边界不清，且有元数据随合成数据传输}
\loc{第202行；第275行；第651行；第692行(v)；对照第188行约束(a)。}
\orig{(第202行) (iii) Synthetic data $\mathcal{D}^{\mathrm{syn}}$ is generated on-node by the local generative model and can be exchanged freely. Therefore, no incomplete or raw metadata leaves a node in unprocessed form.\\
(第275行) Such strong noising departs far from the source volume, so it also favors privacy-preserving de-identification [leiblStablePrivacy2024].\\
(第651行) ... while the metadata text attached to it adds less than 1~KB.}
\ptype{7 事实、推理与证据边界；3 结论放大；4 逻辑跳跃。}
\basis{(a) “can be exchanged freely”是没有证据的确定性断言：本文没有做成员推断、重识别或记忆化检验，第692行局限(v)本身也承认需要额外的隐私措施，前后表述矛盾。(b) 所引 leiblStablePrivacy2024 在 bib 中的题名为 \emph{StablePrivacy: Diffusion-Based Privacy Enhancement for Face Image Datasets}，研究对象是人脸图像；能否支持 3D 脑 MRI 的去识别\verify。(c) 第651行显示，每个合成体积都附带元数据文本，而 img2img 合成以真实受试者的扫描为源，其条件文本很可能来自该受试者经补全的元数据（年龄、性别、健康状态等）。若如此，与约束(a)“no raw $m_j$ may be transmitted”之间的关系需要说明：“补全过”不等于“去识别”。(d) “Therefore”一句与前文“三类数据”划分之间没有推理关系。}
\conf{高（(a)(b)(d)）；中（(c)，需与代码核实条件文本的来源）}
\normal{否。隐私是本文三项核心主张之一，未经评估的确定性表述属于实质性问题。}
\advice{改写；核实合成数据所附元数据的来源与内容；建议补充隐私评估或明确将其列为局限。}
\revise{(iii) Synthetic data $\mathcal{D}^{\mathrm{syn}}$ are generated on-node by the local generative model and are the only data exchanged between nodes. Each synthetic volume is accompanied by its conditioning text \authfill{说明该文本是否来自源受试者的真实/补全元数据}. We use a high noise strength in the federated loop to reduce the correspondence between a synthetic volume and its source; we did not quantify the residual privacy risk in this study (see Limitations).}
\end{finding}

%---------------------------------------------------------------------
\begin{finding}{“可插拔框架”属于未验证的设计主张}
\loc{第206行。}
\orig{It is worth noting that M4Fed is a pluggable framework rather than a fixed, immutable model. The imputation model, the generative model, and the domain model can each be replaced with alternative models ..., so as to accommodate the differing needs of downstream applications and to keep pace with the rapid evolution of foundation models.}
\ptype{3 创新性放大；1 套话（“It is worth noting that”）。}
\basis{本文没有任何替换组件的实验；“keep pace with the rapid evolution of foundation models”是愿景性表述。作为设计说明可以保留，但需标明未验证。}
\conf{中}
\normal{方法论文常说明模块可替换，属正常表达；问题在于语气偏强，且没有说明未经验证。}
\advice{精简。}
\revise{The three models are modular: in principle each can be replaced by an alternative model for a different sub-domain. We did not evaluate such replacements in this study.}
\end{finding}

%---------------------------------------------------------------------
\begin{finding}{TabPFN 补全：动机推理未验证，补全质量与上下文范围未报告}
\loc{第258--262行。}
\orig{Consequently, the conditioning distribution becomes biased toward sub-datasets that happen to report that field.\\
A TabPFN regressor handles continuous fields such as age, weight, and body-mass index ...\\
As a result, this imputation step can be applied independently and consistently at every node without institution-specific calibration.}
\ptype{7 证据边界；8 方法细节缺失；9 一致性。}
\basis{(a) “Consequently ... biased”是合理推断，但没有数据说明（如各字段缺失率按子数据集的分布）。(b) 第345行说明队列只保留年龄非缺失的扫描，因此“age”作为补全对象的例子与数据设置不符，除非补全用于其他场景。(c) 没有报告补全准确性（例如对已观测值做掩码后的补全误差），“consistently ... without institution-specific calibration”无法评估。(d) 没有说明 TabPFN 的拟合上下文是否包含测试划分的元数据行，这关系到是否存在信息泄漏。}
\conf{中}
\normal{动机段落的推断性写法属正常；缺失的补全质量与上下文范围则是复现所需信息。}
\advice{核实并补充。}
\revise{\authfill{各字段缺失率（按节点/子数据集）}；\authfill{补全准确性评估方法与结果}；\authfill{TabPFN 上下文是否包含测试划分}；将 “such as age, weight” 改为实际被补全的字段。}
\end{finding}

%---------------------------------------------------------------------
\begin{finding}{MedMD 段落过长，噪声强度的理由存在推理跳跃}
\loc{第275行。}
\orig{(i) during the federated loop, we use $\sigma > 0.7$ to maximize the diversity of the exchanged synthetic corpus, and thereby the information it carries about the target node's distribution. ... (ii) During local training, we instead sample $\sigma \in [0.35, 0.70]$ ... In this range, the augmented volume deviates from the original but still retains part of its class information [sastryDiffAug2024].}
\ptype{4 逻辑跳跃；5 超长段落；7 证据边界。}
\basis{(a) 噪声越强，合成体积与源体积的对应越弱；“多样性最大化”并不意味着“携带更多关于目标节点分布的信息”，两者甚至可能此消彼长，这一推理需要论证或改为假设。“maximize”与固定取值 $\sigma=0.75$（表4）也不相符，因为并未做优化。(b) DiffAug 研究的是自然图像分类，“$[0.35,0.70]$ 保留类别信息”能否迁移到脑龄回归\verify，且与第645行“CV 领域的范围不能直接迁移”的结论构成张力（见条目34）。(c) 本段约 350 词，混合了模块作用、结构、采样流程、噪声日程与训练目标，建议拆分。(d) “Given a target conditioning string”没有说明目标文本是源受试者自身的元数据还是经修改的元数据。}
\conf{中--高}
\normal{设计理由段落常带有推测成分，但应标为设计假设。}
\advice{改写并拆段。}
\revise{During the federated loop we use a high noise strength ($\sigma=0.75$), so that the exchanged volumes depart further from their sources; we expect this to increase the diversity of the exchanged data, at the cost of weaker correspondence to individual source scans. During local training we sample $\sigma\in[0.35,0.70]$ uniformly, a range reported to preserve class information for natural-image classifiers [sastryDiffAug2024]; whether it transfers to brain-age regression is examined in Section~5.7. \authfill{说明条件文本的来源}}
\end{finding}

%---------------------------------------------------------------------
\begin{finding}{方法与实验设置中缺少复现所需的步骤、参数和统计单位}
\loc{第4.4、5.1、5.2节；表4；第298行；第365行；图 test\_scatter\_registered。}
\orig{(第298行) Samples with a missing label for field $q$ are masked out of $\mathcal{L}_q$ rather than imputed, since imputation here would contaminate the evaluation signal ...\\
(第365行) Within each private split, $90\%$ of samples are used for training and $10\%$ for held-out evaluation.}
\ptype{8 方法缺失步骤、参数和统计单位。}
\basis{以下信息缺失或含糊：(1) 每轮本地训练的 epoch 数、优化器、学习率调度；(2) DDIM 的 $\eta$、是否使用 classifier-free guidance 及其系数；(3) MedMD U-Net 上 QLoRA 的秩和注入层；(4) 基线的通信轮数（表4中 $R_f=3$ 是否也是 FedAvg 的轮数，3 轮能否使基线收敛）；(5) 服务器下发的合成体积数（$\lceil\rho_f n_s^{\mathrm{priv}}\rceil$ 中 $n_s^{\mathrm{priv}}$ 未给出）；(6) 随机种子与重复次数，“mean $\pm$ std”的统计单位（M4Fed 为 2 个客户端，基线为 3 个）；(7) 90/10 划分的单位：OpenMind 中 26{,}408 次扫描对应 20{,}107 名受试者，存在同一受试者多次扫描的情况。健康对照队列每块约 1{,}143 名受试者，但图 test\_scatter\_registered 中各节点测试集为 169--214 个样本，明显超过 10\% 的受试者数，说明测试单位可能是扫描；若按扫描划分，同一受试者可能同时出现在训练集和测试集中。(8) 第298行“contaminate the evaluation signal”：被掩码的是训练标签，措辞应改为“supervision signal”，或说明评估时如何处理缺失标签。}
\conf{高（第7点需核实；其余为可直接观察到的缺失）}
\normal{否，属于复现必需信息。}
\advice{核实并补充（可放补充材料）。}
\revise{Within each private block, subjects (not scans) are split 90\%/10\% into training and test sets, so that no subject contributes scans to both \authfill{核实实际划分单位}. All results are from a single run with seed \authfill{种子} \authfill{或报告 N 次重复的均值与标准差}.}
\end{finding}

%---------------------------------------------------------------------
\begin{finding}{复杂度分析中“decoupled”“solely”与同段内容矛盾；$|\phi|$ 定义混乱}
\loc{第313、315行；表2。}
\orig{Consequently M4Fed's communication cost is mainly $\rho_f\, n_k\, \beta$, metered solely by the synthesis ratio $\rho_f$ and decoupled from the size of the domain and generative backbones ...\\
let $|\phi| = r(d_{\mathrm{in}}+d_{\mathrm{out}})$ denote the per-layer QLoRA adapter size at rank $r$, summed over adapted layers, plus the task-head parameters, so that $|\phi| \ll |\theta_{\mathrm{DM}}|, |\theta_{\mathrm{GM}}|$ by construction.}
\ptype{9 一致性；4 概念含糊；3 放大。}
\basis{(a) 同一段落中已说明一次性广播 $|\theta_{\mathrm{GM}}|$，它与生成骨干的大小相关；单轮开销还取决于 $n_k$ 和单体积大小 $\beta$（即分辨率），因此“solely”“decoupled from ... generative backbones”不准确。表2也没有列出服务器→客户端每轮下发的 $\lceil\rho_f n_s\rceil\beta$。(b) $|\phi|$ 先定义为“per-layer”，随后又称“summed over adapted layers, plus task heads”，一个符号承担了三种含义。(c) “$\ll$ by construction”：实测 $|\phi|=1.67$M，约为 $|\theta_{\mathrm{DM}}|$ 的 18\%，其中任务头 1.19M 未被压缩，“$\ll$”的强度偏高。}
\conf{中--高}
\normal{符号化复杂度分析常做简化，但这里的简化与同段表述冲突。}
\advice{改写。}
\revise{Per round, each client uploads $\lceil\rho_f n_k\rceil\beta$ bytes and receives $\lceil\rho_f n_s^{\mathrm{priv}}\rceil\beta$ bytes of synthetic data; in addition, the trainable part of the generative model ($|\theta_{\mathrm{GM}}|$) is broadcast once after warm-up. The per-round payload therefore depends on $\rho_f$, the node sample sizes, and the volume size, but not on $|\theta_{\mathrm{DM}}|$. Let $|\phi| = \sum_{\ell} r\,(d^{\ell}_{\mathrm{in}}+d^{\ell}_{\mathrm{out}}) + |\theta_{\mathrm{head}}|$ denote the number of client-side trainable parameters.}
\end{finding}

%---------------------------------------------------------------------
\begin{finding}{数据集段落：语法、空泛句与未在正文说明的排除步骤}
\loc{第345、347行；表3说明（第351行）。}
\orig{Filtering to entries with a non-missing age label, retains $26{,}408$ T1w scans ...\\
... and lifespan brain-age prediction has become a potential research direction [wu2024lifespan]. In addition, restricting the data to adults would also discard the many pediatric OpenNeuro sites ...\\
(表3) after excluding the $54$ corrupted and $1$ non-finite-voxel scans ... the deduplicated combine is their deduplicated union.}
\ptype{5 语法；1 空泛；9 一致性。}
\basis{第一句主语后多一个逗号。“has become a potential research direction”信息量很低。“the many pediatric OpenNeuro sites”没有给出数量。损坏扫描的排除只出现在表注中，正文流程（$42{,}732\to26{,}408$）里没有体现。“deduplicated combine”应为“combined”。“approximately 800 OpenNeuro sub-datasets”与“$42{,}732$ of the $113{,}921$”需要与 OpenMind 原文核对（\verify）。}
\conf{低--中}
\normal{多数为语言层面的小问题。}
\advice{精简并补充。}
\revise{Restricting the data to entries with a non-missing age label leaves $26{,}408$ T1w scans from $20{,}107$ subjects; we further exclude $54$ corrupted scans and $1$ scan with non-finite voxels. ... Both cohorts include minors (4--18 yr), because lifespan models can reveal disorder-specific brain-age patterns [kaufmann2019common], and excluding minors would remove \authfill{数量} pediatric sub-datasets.}
\end{finding}

%---------------------------------------------------------------------
\begin{finding}{实验设置中节点与数据块的对应关系含糊；基线描述与补充材料不一致}
\loc{第367、371行；补充材料 S 第176、181行；第156行。}
\orig{FedAvg, FedProx, FAN, and GPT-FL train on three clients with the same private block, whose server only performs aggregation and broadcasting while GPT-FL's server reuses the generator trained by M4Fed's.\\
(第371行) These baselines represent increasing degrees of training-resource utilization: ... FAN additionally exploits synthetic data ...\\
(S 第176行) FAN is a decentralized baseline ... in which three clients hold independently initialized MedMD and SwinMTL instances and perform no server-side aggregation.}
\ptype{9 一致性；5 表达不清；8 方法描述。}
\basis{(a) “with the same private block”含义不明：是三客户端分别使用 M4Fed 的三个私有块，还是共用同一块？(b) 第367行说基线的服务器“performs aggregation and broadcasting”，而补充材料说 FAN 没有服务器端聚合。(c) 第156行把 FAN 称为“Federated Adversarial Network”，补充材料中的复现却基于 MedMD 扩散 img2img；GPT-FL 原方法用 LLM 生成提示词，复现中改为“conditioned on age label seed alone”。复现与原方法的差异应在正文中说明，否则“re-implemented ... to isolate the effect of the federated protocol”的公平性难以评估。(d) 第371行把“synthetic data”当作一种资源，与表1中“三类资源”的口径不同。(e) “all experiments use a NVIDIA A100 GPU”应为“an”。}
\conf{高（(a)(b)）；中（(c)(d)）}
\normal{否，影响读者对比较公平性的判断。}
\advice{改写，并在补充材料中列出复现与原方法的差异。}
\revise{For each baseline, the three private blocks are assigned to three clients, \authfill{核实：是否与 M4Fed 的服务器块和两个客户端块一一对应}. FedAvg, FedProx, and GPT-FL use a server that only aggregates and broadcasts weights; FAN has no server. Our re-implementations differ from the original methods as follows: \authfill{FAN 与 GPT-FL 的改动}.}
\end{finding}

%---------------------------------------------------------------------
\begin{finding}{主比较的比较单位不对等，“by a clear margin”缺少统计支撑}
\loc{第448行；表5。}
\orig{M4Fed server is the strongest model by a clear margin on every metric ... M4Fed clients point-error metrics (MAE $10.97$, RMSE $17.19$) are not the best in the table, being comparable to or slightly worse than FedAvg, FedProx, and GPT-FL. ... so the server's accuracy advantage is expected.}
\ptype{3 结论放大；9 一致性；8 统计单位。}
\basis{(a) M4Fed 服务器是单个节点的单次结果，同时拥有公开数据和全参数微调；基线是 3 个客户端的均值。两者的测试集、训练资源与统计单位都不同，文中也承认这种优势“is expected”。若主结论（摘要中的“surpasses”）主要依赖这一对比，需要说明其性质。(b) 没有重复运行和显著性检验，客户端标准差约为均值的 40--50\%，“clear margin”缺少统计依据。(c) M4Fed 客户端的 RMSE 是全表最高，$R^2$ 是全表最低，“comparable to or slightly worse”的说法偏轻。(d) “under the identical environment”含义不明。}
\conf{高}
\normal{客户端/服务器分开报告本身合理；问题在于结论性措辞没有反映比较条件。}
\advice{改写；建议补充按节点逐一对应的比较（例如基线在服务器私有块测试集上的结果）和重复实验。}
\revise{The M4Fed server attains the lowest MAE (8.87 vs.\ 10.28 years for GPT-FL) and is the only model with a positive $R^2$. This comparison is between a single server node, which additionally has public data and full-parameter fine-tuning, and the client average of each baseline; \authfill{给出同一测试块上的基线结果或说明无法对齐的原因}. The M4Fed clients have higher MAE and RMSE and the lowest $R^2$ of all methods, but the highest Pearson correlation (0.442).}
\end{finding}

%---------------------------------------------------------------------
\begin{finding}{把 Pearson $r$ 称为“排序”指标}
\loc{第455行；第684行；对照第448行。}
\orig{(第455行) ... suggesting that these two components matter more for the ranking quality of predictions than for their point-wise calibration. ... This result indicates that a stronger rank correlation coexists with worse absolute calibration.\\
(第684行) ... they still rank subjects by age more accurately than every baseline on Pearson $r$ ...}
\ptype{4 术语含义不准确。}
\basis{Pearson $r$ 衡量的是线性相关，不是秩相关（秩相关对应 Spearman $\rho$，第677行已正确使用）。“ranking quality”“rank correlation”“rank subjects”与所报告的指标不符。另外，Pearson 高而 $R^2$ 为负，说明预测存在尺度或偏移偏差，这与图 test\_scatter\_registered 中预测值集中在约 23--28 岁的现象一致，可以在文中直接说明。}
\conf{高}
\normal{否，属于术语误用。}
\advice{改写；或改为同时报告 Spearman $\rho$。}
\revise{... suggesting that these components affect the linear association between predicted and true age more than the absolute error. ... a higher linear correlation coexists with worse absolute calibration.\\
(第684行) ... their predictions are more strongly correlated with chronological age than those of every baseline (Pearson $r$ 0.442 vs.\ 0.248) ...}
\end{finding}

%---------------------------------------------------------------------
\begin{finding}{组件消融：“confirming”结论强于证据，三个消融结果几乎相同，需要核实}
\loc{第455行；表6（第458--472行）；第684、686行。}
\orig{Every ablation increases MAE relative to the full model, by a fairly uniform $5.4$--$5.6\%$ ..., confirming that public-data warm-up, TabPFN-densified conditioning, and pretrained foundation model each contribute positively. ... Removing public data yields the worse $R^2$ ($-0.264$) and MAE ($11.56$) ...}
\ptype{3 结论放大；9 一致性；7 证据边界。}
\basis{(a) 三个消融变体的 MAE（11.56/11.58/11.57）、MAE 标准差（4.98/4.98/4.98）、RMSE（18.79/18.84/18.81）和 $R^2$（$-0.264/-0.272/-0.268$，标准差 0.000--0.002）几乎完全相同。三个性质完全不同的改动（去公开数据、去补全、随机初始化）产生几乎一致的结果，更像是三者都退化成了相近的近常数预测器。这需要核实，在核实之前不宜把差异分别归因于各组件。(b) 差异约 0.6 年，而客户端间标准差约 5 年，又是 2 个客户端的单次运行，“confirming”过强。(c) “the worse $R^2$ ($-0.264$)”：在三个消融中 $-0.264$ 其实是最好的，措辞容易误读。(d) 消融只在客户端上报告，而全文的主要精度结论来自服务器。(e) 第686行称消融“verifies that MedMD can enhance the performance of the domain model”，但没有任何消融去掉 MedMD 或合成数据。(f) 随机初始化骨干并冻结后仅训练 QLoRA 适配器，这一设置的含义需要说明。}
\conf{高（(c)(d)(e)）；中（(a)，需核实原始预测）}
\normal{对推测性解释使用“may”“possible”是恰当的；但“confirming”与数据强度不匹配。}
\advice{核实（检查三个消融变体的预测分布）；改写结论措辞；补充服务器端消融与去除合成数据的消融。}
\revise{All three ablations increase client-average MAE by about 0.6 years (5.4--5.6\%), a difference that is small relative to the variation between clients (std $\approx$ 5 years) and is based on a single run; we therefore read it as an indication rather than confirmation that each component contributes. \authfill{三个变体预测分布的核查结果}}
\end{finding}

%---------------------------------------------------------------------
\begin{finding}{同一默认配置在三张表中给出三组不同的服务器结果}
\loc{表5（第442行）；表8（第528、533行）；表7（第493行）；第517、690行。}
\orig{(表5) M4Fed (server): MAE 8.87, RMSE 14.36, $r$ 0.554, $R^2$ 0.178\\
(表8, $\rho_f=0.1$ 与 MedMD lr $=10^{-4}$，即表4默认值) MAE 9.376, RMSE 15.172, $r$ 0.417, $R^2$ 0.061\\
(表7, 未筛选队列, Mixed/Mixed) MAE 9.22, RMSE 15.11, $r$ 0.421, $R^2$ 0.068}
\ptype{9 内容一致性。}
\basis{按表4的默认超参数，表8中这两行应与表5的服务器行为同一配置，但数值不同；表7的 Mixed/Mixed 行也应是同一设置。第517、690行又以 8.87 作为“default partition”的参照。如果差异来自不同的运行（随机种子），那么同一配置下的 MAE 波动约 0.5 年、Pearson 在 0.42--0.55 之间，与消融效应（约 0.6 年）和敏感性分析中的差异量级相当，许多“robust”“contribute”的结论都需要重新表述；如果来自不同设置，则需要说明区别。}
\conf{高}
\normal{否。}
\advice{核实并在表注中说明各表所用的运行与设置。}
\revise{\authfill{说明表5、表7、表8的服务器结果是否同一配置；若为不同运行，报告多次运行的均值与标准差，并据此修改第5.2--5.5节中所有基于单次差异的结论}}
\end{finding}

%---------------------------------------------------------------------
\begin{finding}{健康状态消融：基于年龄方差极小的子集给出部署建议，段内逻辑顺序跳跃}
\loc{第477--482行；表7。}
\orig{... cutting healthy-test MAE by $74\%$ ($1.66$ vs.\ $6.41$) ... This result seems to indicate that when predicting the age of patients with unknown health conditions, mixed data should be used for training, while when predicting the age of healthy individuals, healthy data should be used for training. It is worth noting that the brain age prediction model was used to assess ... Therefore, the main focus should be on the results where the test slice is healthy people. For this part, we flag statistical caution ...\\
(第482行) In order to resolve this doubt, we therefore add ... This value is slightly lower than the mixed-train, healthy-test MAE on the unselected cohort ($6.41$), but the RMSE metric (up to $12.45$) does not show a corresponding gain. Therefore, restricting training to healthy subjects does not necessarily yield a model that explains more age variance on the healthy test slice.}
\ptype{7 证据边界；4 逻辑跳跃；3 放大；5 冗余。}
\basis{(a) 由 $\mathrm{Var}(y)=\mathrm{RMSE}^2/(1-R^2)$ 可从表7反推健康测试切片的年龄标准差：$6.50/\sqrt{40.58}\approx1.0$ 年，$1.95/\sqrt{3.667}\approx1.0$ 年，两行一致。也就是说，该切片的年龄几乎相同，MAE 1.66 很可能只反映了预测值接近该切片的均值，不足以支持“specializes the model”或“healthy data should be used for training”这一部署建议。文中随后也提出了谨慎，但给建议在前、提醒在后，顺序颠倒。各测试切片的样本量没有报告。(b) 段落推进顺序是：建议 → 脑龄用途 → “Therefore”关注健康组 → 统计谨慎，各句之间的因果连接词（“Therefore”）不成立。(c) 第482行把健康对照队列的 MAE（6.18）与未筛选队列健康切片的 MAE（6.41）比较，二者测试集和年龄分布完全不同，不可直接比较；最后的“Therefore”结论基于 RMSE 推出“explains more age variance”，指标与结论不对应。(d) “In order to resolve this doubt, we therefore”冗余。}
\conf{高}
\normal{否；(a) 是可由稿件自身数字推出的证据问题。}
\advice{改写；补充各切片样本量与年龄分布。}
\revise{Training on healthy subjects only reduces the MAE on the healthy test slice (1.66 vs.\ 6.41 years) but increases the MAE on the mixed test set by 21\% (11.15 vs.\ 9.22). The healthy test slice is, however, small (\authfill{n}) and its age spread is very narrow (an age SD of about 1 year, implied by the reported RMSE and $R^2$), so its MAE mainly reflects how close the predictions are to the slice mean, and neither $r$ nor $R^2$ is informative there. We therefore do not draw conclusions about training on healthy subjects from this slice. To obtain a healthy test set with a realistic age range, we additionally train and test M4Fed on the healthy-controls cohort (Table~7, last row: MAE 6.18 years, $R^2=-0.045$); this result is not directly comparable to the unselected-cohort rows because the test sets differ.}
\end{finding}

%---------------------------------------------------------------------
\begin{finding}{敏感性分析：“robust”“adapt well”与数据不完全一致，解释多为未验证的推测}
\loc{第515、517行；表8。}
\orig{This indicates that M4Fed can maintain performance when significantly reducing the amount of synthetic data ..., so that it has potential in terms of largely reducing communication burden without performance loss.\\
... showing that M4Fed can adapt well to heterogeneous data distributions across nodes. We see two plausible reasons for this: (i) ... (ii) ...\\
A possible explanation is that a higher degree of non-IID skew makes the data within each node more long-tailed, and M4Fed's local data augmentation reinforces this tendency ... This makes it easier for the model to learn the overall trend of the age-group distribution ...}
\ptype{3 结论放大；7 推测写成解释；2 句式重复。}
\basis{(a) 每个轴只有 3 个取值、单次运行，没有重复实验（见条目27），“robust”“without performance loss”应限定为“在所测范围内”。(b) Dirichlet 重划分后 MAE 比默认划分高约 20\%，$R^2$ 全部为负，$\alpha=1.0$ 时 Pearson 为 $-0.195$，与“adapt well”不符；同段后半句又承认“has limits”，前后语气不一致。(c) 两条原因（预训练表征、合成数据增强）没有对应实验检验，例如在偏斜划分下去掉预训练或合成数据。(d) Pearson 的解释链条为“偏斜 → 长尾 → 增强加剧 → 更易学到总体趋势 → 相关性提高”，每一步都未经验证；$r=-0.195$ 这一反常值也可能来自训练失败，需要先排查。(e) 重划分打破了子数据集边界，与默认的按子数据集划分不是同一种数据分布，与 8.87 直接比较会混入划分方式的影响。(f) “We see two plausible reasons”与第574行句式完全相同。}
\conf{中--高}
\normal{对推测性解释使用“possible”是恰当的；问题在于“adapt well”等结论性措辞，以及推测链过长。}
\advice{改写；核实 $\alpha=1.0$ 的运行。}
\revise{Within the tested range, server MAE varies by at most 5.3\% across $\rho_f\in\{0.01,0.1,1.0\}$ and MedMD learning rates $\{10^{-6},10^{-5},10^{-4}\}$ (single runs). Under Dirichlet re-partitioning, server MAE varies by 2.0\% across $\alpha$, but is about 20\% higher than under the default sub-dataset partition and $R^2$ becomes negative; because re-partitioning also mixes sub-datasets across nodes, this difference cannot be attributed to quantity skew alone. The negative Pearson $r$ at $\alpha=1.0$ \authfill{核实是否为训练异常}.}
\end{finding}

%---------------------------------------------------------------------
\begin{finding}{配准实验：预测值近乎常数时，用 MAE 变化说明域模型改善需要加以限定}
\loc{第574行；表 registration；图 test\_scatter\_registered（图 age-zscore-scatter (a)）。}
\orig{For the domain model, the effect is smaller on clients, but is still visible on the server: MAE drops by $4.2\%$ ... In summary, affine registration is a low-cost preprocessing step that can further improve generative model and server's domain model. In addition, registration requires only a public template and is performed locally before training, it does not affect ... our results support the conclusion of [rajabli2025brainage] ...}
\ptype{7 证据边界；3 放大；5 语法。}
\basis{(a) 图 test\_scatter\_registered 显示，配准后（未标准化）各节点的预测值几乎都集中在约 23--28 岁，与真实年龄（6--79 岁）基本无关。在这种近常数预测下，MAE 从 6.18 降到 5.92，主要反映常数值与 20--30 岁密集人群的距离，而不是年龄建模能力的提升，应在文中说明。(b) 原因 (i) 中“latent standard deviation 0.82→0.89”的数据只出现在这里，未在表中报告；原因 (ii)（“already partly invariant to position”）是推测。(c) “low-cost”没有测量配准耗时。(d) 最后一句是逗号连接两个独立句。(e) 单队列、单次运行，“support the conclusion of”可改为“is consistent with”。(f) 这里 VAE 被描述为“pretrained on co-registered BraTS volumes [antonelli2022msd]”，补充材料 S 第137行则写作“pretrained for brain MRI generative modeling”，两处应统一并给出具体 checkpoint。}
\conf{高（(a)）；中（其余）}
\normal{(b)(e) 的语气稍作调整即可；(a) 属实质性问题。}
\advice{改写。}
\revise{For the domain model, registration lowers server MAE by 4.2\% (6.18 to 5.92 years) and raises Pearson $r$ from 0.249 to 0.381. Without label standardization, however, the predictions of all nodes remain concentrated around 23--28 years (Fig.~4a), so these changes mainly reflect a shift of the near-constant prediction rather than better age modeling. ... Because registration uses only a public template and is performed locally before training, it does not change the privacy or communication properties of the protocol. These results are consistent with [rajabli2025brainage], who report that registration remains beneficial for brain-age models.}
\end{finding}

%---------------------------------------------------------------------
\begin{finding}{标签标准化：缺少谓语的句子；所给机制未经验证，可能不成立}
\loc{第593行。}
\orig{Figure~4 may show where these gains come mainly from. Without standardization, the predictions of every node mainly around the young-adult majority, while the standard deviation of the predicted ages increases with standardization. ... Our results suggest that label standardization can partly mitigate this problem because standardization can alleviate the situation where the scale of age labels leads to the model giving more weight to dense samples and neglecting sparse samples.}
\ptype{5 语法；7 推测写成因果；4 逻辑跳跃。}
\basis{(a) 第二句缺少谓语（应为“are concentrated”）。“may show where these gains come mainly from”用语含糊。(b) z-score 是仿射变换；在 MSE 损失下，它只按常数 $1/\sigma_k^2$ 缩放损失，并不改变密集样本与稀疏样本之间的相对权重。所述机制（“scale of age labels leads to the model giving more weight to dense samples”）既未经实验验证，也不能从变换本身推出。更可能的作用途径（例如有效梯度尺度、输出层初始化与目标量级的匹配）同样未验证，因此应写为假设，或不给机制。(c) 服务器 Pearson 从 0.381 降到 0.345，结果并非一致改善；第690行的总结（“registration and label standardization can improve”）没有反映这一点。}
\conf{中--高}
\normal{(a) 为语言问题；(b) 属于把推测写成因果。}
\advice{改写。}
\revise{Figure~4 compares the predictions with and without standardization. Without standardization, the predictions of all nodes are concentrated around the young-adult majority; with standardization, their spread increases. This regression towards the mean, also known as age-dependent bias, is commonly reported in brain-age studies [smith2019estimation, liang2019investigating] and is associated with sparse samples at the extremes of the age range [yang2021delving]. Standardization partly reduces this effect for the clients in our experiments; we have not identified the mechanism.}
\end{finding}

%---------------------------------------------------------------------
\begin{finding}{跨站点泛化：负 $R^2$ 被解释为“并非模型失效”，跨测试集比较被称为“conservative”}
\loc{第602、604行。}
\orig{Hence, the drop of $r$ and $R^2$ mainly reflects the composition of $\mathcal{D}^{\mathrm{cross}}$ rather than a failure of the model ...\\
Judged by these point-error metrics, M4Fed generalizes well to unseen sites. Its cross-site MAE ($9.42$) remains lower than the in-distribution MAE of every baseline ... This comparison is conservative, because ... We attribute this robustness to two parts of the design.}
\ptype{3 结论放大；7 证据边界；4 逻辑跳跃。}
\basis{(a) 文中关于“方差变窄导致 $r$、$R^2$ 下降”的数值解释是自洽的（由 RMSE 与 $R^2$ 反推的标准差为 13.3 年，与文中一致）。但 $R^2<0$ 意味着模型不如直接预测跨站点集合的均值；该集合约 70\% 的扫描年龄在 18--24 岁，9.42 年的 MAE 相对这一年龄带仍然很大。“rather than a failure of the model”是超出证据的判断。(b) 拿 M4Fed 在 $\mathcal{D}^{\mathrm{cross}}$ 上的 MAE 与基线在各自分布内测试集上的 MAE 比较，测试集不同，年龄分布也不同。同段前文刚论证过指标依赖评估集的年龄分布，而 MAE 同样依赖它，因此“conservative”无法成立。基线没有在 $\mathcal{D}^{\mathrm{cross}}$ 上评估。(c) “We attribute this robustness to two parts of the design”：两条归因都没有在 $\mathcal{D}^{\mathrm{cross}}$ 上做消融检验，所引的无公开数据消融是客户端的分布内结果。(d) 约四分之三的跨站点扫描为患者，这与该池“健康状态缺失率 <50\%”的筛选条件之间的关系需要说明。}
\conf{高}
\normal{否。}
\advice{改写；建议在 $\mathcal{D}^{\mathrm{cross}}$ 上评估所有基线。}
\revise{On $\mathcal{D}^{\mathrm{cross}}$, the server's MAE (9.42 years) and RMSE (13.99 years) are close to its in-distribution values, whereas Pearson $r$ drops to 0.202 and $R^2$ becomes negative. Part of this drop is expected from the narrower age distribution of $\mathcal{D}^{\mathrm{cross}}$; a negative $R^2$ nevertheless means that the model does not outperform a constant prediction at the cross-site mean. \authfill{基线在 $\mathcal{D}^{\mathrm{cross}}$ 上的结果}. Public-data warm-up and synthetic exchange may contribute to the stable point errors, but we did not test this on $\mathcal{D}^{\mathrm{cross}}$.}
\end{finding}

%---------------------------------------------------------------------
\begin{finding}{生成质量：服务器较差的解释只能说明“不更好”，不能说明“更差”；评估设置未说明}
\loc{第643行；表9。}
\orig{Interestingly, M4Fed server row is weaker than its own client-average row on LPIPS/PSNR/KID. A possible reason is that the server's larger data pool does not actually benefit its generative model. ... Hence, compared with the clients, the server gains no extra real-image supervision for MedMD, and its data advantage benefits the domain model rather than synthesis quality.}
\ptype{4 逻辑跳跃；7 证据边界；8 方法细节。}
\basis{(a) “gains no extra real-image supervision”最多能解释服务器为何不比客户端好，不能解释为何更差。(b) “its data advantage benefits the domain model rather than synthesis quality”未经检验。(c) 表9没有说明评估所用的队列和噪声强度 $\sigma$；PSNR、SSIM 是合成体积与其源体积配对计算的，强烈依赖 $\sigma$，不同方法若 $\sigma$ 不同则不可比。(d) FID 约 300 表示与真实分布差距很大，且使用的是自然图像预训练的 Inception 网络，解读时应注明。(e) “Interestingly”在第455行也出现过，属于模板化过渡。}
\conf{中}
\normal{对意外结果给出可能原因属正常做法；但第二句起使用了确定语气。}
\advice{改写并补充设置。}
\revise{The server row is weaker than the client average on PSNR, LPIPS, and KID. Because all nodes start the federated loop from the same warm-up weights, the server receives no additional real-image supervision for MedMD; this explains why it is not better than the clients, but not why it is worse, which we leave open. \authfill{表9所用队列与 $\sigma$}}
\end{finding}

%---------------------------------------------------------------------
\begin{finding}{噪声强度实验的结论超出证据：“CV 领域的范围不能迁移到医学影像”}
\loc{第645行；图 qualviz。}
\orig{Lowering $\sigma$ turns the blurry synthesized slices into slices with clear anatomical structure. We therefore ran an additional hyperparameter analysis ... using the healthy-subjects-only OpenMind subset. ... Taken together, these results suggest that the range of suitable noise strengths established in the computer vision domain does not transfer directly to medical imaging. Instead, the appropriate noise-strength range should be explored for the specific imaging modality and data under study.}
\ptype{3 局部结果推广；4 逻辑跳跃；9 术语一致性。}
\basis{(a) 实验只比较了默认设置（联邦循环 0.75，本地训练 $[0.35,0.70]$）与 $\sigma=0.05$，结果是下游 MAE 几乎不变。这既不能说明 CV 领域的范围不适用，也不能说明它适用。从一个数据集、两个设置、单次运行推广到“medical imaging”，范围过大。(b) $\sigma=0.05$ 时输出接近源图像的重建，切片更清晰是预期结果，作为发现来陈述的信息量有限。(c) “healthy-subjects-only OpenMind subset”与全文的“healthy-controls cohort”术语不一致。(d) 段内的“parallel”结论（“does not translate into a better domain model”）是合理的。}
\conf{高}
\normal{否。}
\advice{改写。}
\revise{... using the healthy-controls cohort. ... Hence, in this setting, a lower noise strength brings the synthetic volumes closer to the real data distribution but does not improve the domain model, possibly because volumes synthesized at $\sigma=0.05$ add little diversity. Whether the noise-strength range used for natural images is appropriate for brain MRI would require a finer sweep of $\sigma$, which we did not perform.}
\end{finding}

%---------------------------------------------------------------------
\begin{finding}{通信与计算实测：只比较单客户端单轮上传；“compute”没有实测}
\loc{第651--653行；第688行。}
\orig{Based on the above basic data, the quantitative improvement effects in terms of computing and communication burdens can be analyzed as follows. ... Moreover, M4Fed's payload does not depend on the size of either backbone ... Moreover, the sensitivity analysis (Table~8) shows that reducing $\rho_f$ to $0.01$ does not degrade server accuracy.}
\ptype{9 一致性；3 放大；1 套话；2 重复连接词。}
\basis{(a) 本节所有数值复算无误。(b) 比较只计入客户端上传：FedAvg 每轮还有 37.0~MB 下行；M4Fed 客户端每轮接收服务器合成数据 $\lceil\rho_f n_s^{\mathrm{priv}}\rceil\beta$（$n_s^{\mathrm{priv}}$ 未给出），另有一次性 30.1~MB 的 U-Net 广播。没有给出 $R_f=3$ 轮的总通信量，也没有说明 FedAvg 达到收敛所需的轮数。(c) “does not depend on the size of either backbone”与一次性广播 $|\theta_{\mathrm{GM}}|$ 相矛盾（见条目21）。(d) 只统计了可训练参数，没有测量显存、训练时间或 FLOPs，而客户端额外的 DDIM 50 步生成开销没有计入，因此“computing burden”的量化并不完整。(e) “does not degrade server accuracy”：$\rho_f=0.01$ 的 MAE 8.900 优于表8的默认行 9.376，但略差于表5的 8.87（见条目27），且是单次运行。(f) 首句是空泛的过渡；“Moreover”连续两句使用。}
\conf{高}
\normal{(a) 是本稿写得较好的部分；(b)(d) 属于口径不完整。}
\advice{补充（双向、总轮次通信量；客户端显存与时间）；精简过渡句。}
\revise{We now instantiate Table~2 with measured values. ... \authfill{双向、$R_f$ 轮总通信量对比表}；\authfill{客户端峰值显存与每轮耗时，含 DDIM 生成}. In the sensitivity analysis, server MAE at $\rho_f=0.01$ (8.90 years) is not higher than at $\rho_f=0.1$ (9.38 years, single runs).}
\end{finding}

%---------------------------------------------------------------------
\begin{finding}{集中式实验：“upper bound”与训练预算不对等，“This gap”指代不明}
\loc{第658、675行。}
\orig{The training data and test data are pooled at a single node respectively, and SwinMTL is fine-tuned for $10$ rounds. ...\\
As shown in Figure~6a, the centralized model reaches an MAE of $3.51$, and the predictions are nearly spread along the identity line although there is still bias over age. This gap is expected, because ... It should therefore be read as an upper bound ...}
\ptype{7 证据边界；9 一致性；5 指代不明。}
\basis{(a) 集中式训练 10 轮，联邦本地训练 $R_l=3$ 轮，训练预算不同，因此不能严格视为“upper bound”。(b) “This gap”没有给出对比对象（应为同设置下 M4Fed 服务器的 MAE，例如配准+标准化后的 5.85），读者需要自己去找。(c) “pooled at a single node respectively”含义不清：公开数据块是否参与训练？(d) 图中 70--80 岁组 MAE 为 23.34 年，“nearly spread along the identity line”与之不太相称。(e) 第677行的统计检验写得规范，可以作为全文的参照。}
\conf{中}
\normal{集中式参照常被称为 upper bound，但前提是训练条件可比。}
\advice{改写。}
\revise{The centralized model reaches an MAE of 3.51 years, compared with \authfill{对应 M4Fed 服务器值} years for the M4Fed server under the same preprocessing. Because the centralized model is trained on all real data with full-parameter updates for 10 rounds (vs.\ 3 local rounds in M4Fed), it serves as a reference for the cost of the federated constraints rather than as a strict upper bound. Predictions follow the identity line for most subjects under 50 years but underestimate older subjects (Fig.~6).}
\end{finding}

%---------------------------------------------------------------------
\begin{finding}{讨论第一、二段：把设计目标写成“主要发现”；跨实验比较不成立；连接词误用}
\loc{第684、686行。}
\orig{(第684行) These results lead to the main finding of this study: M4Fed provides a multi-source federated learning framework for jointly and effectively exploiting the training resources ...\\
(第686行) On the generative side, the component ablation verifies that MedMD can enhance the performance of the domain model by leveraging the imputed multi-modal data. ... On the contrary, this might be because the Swin UNETR encoder ... can better learn position invariance ... Therefore, the impact of registration preprocessing and changing noise intensity is limited, while imputed multi-modal data are more influential.}
\ptype{6 信息增量不足；3 放大；4 逻辑跳跃；5 连接词误用。}
\basis{(a) “main finding”是对方法设计目标的复述，不是实验发现；“effectively”与客户端结果不符（条目24）。(b) 没有任何消融去掉 MedMD（条目26），“verifies that MedMD can enhance”缺少依据。(c) “imputed multi-modal data are more influential”比较的是未筛选队列上客户端的补全消融（MAE +5.6\%）和健康队列上服务器的配准效果（MAE $-4.2\%$），队列、节点和方向都不同，无法得出“更有影响力”。(d) “On the contrary”与前句并非对立关系；“this might be because”之后紧接“Therefore”给出确定结论，推测被当作了前提。(e) 第686行称配准影响“limited”，第690行又称“registration ... can improve”，两处口径不一致。}
\conf{高}
\normal{否。}
\advice{改写。}
\revise{The M4Fed server achieves the lowest MAE of all methods (13.7\% below the best baseline), whereas the clients improve only the correlation-based metric. The component ablation indicates small contributions of public data, imputation, and pretraining at the clients (Section~5.3). Synthesis quality and downstream accuracy appear only loosely coupled in our experiments: lowering $\sigma$ or registering the inputs improves FID substantially while changing server MAE by at most 4.2\%. One possible reason is that the Swin UNETR encoder is partly invariant to spatial variation; we did not test this.}
\end{finding}

%---------------------------------------------------------------------
\begin{finding}{讨论第三、四段：主要是结果复述，归因未经检验}
\loc{第688、690行。}
\orig{(第690行) These results suggest that pretrained representations and synthetic exchange give M4Fed a reasonable generalization under distribution shift, but do not fully preserve the finer label ordering learned within the federation. With respect to preprocessing, the experiment found that registration and label standardization can improve the brain age estimation based on M4Fed.}
\ptype{6 重复与信息增量不足；7 证据边界。}
\basis{第688行基本重复第5.8节的数字（81.9\%、92.6\%、5.55$\times$、0.018、19.5$\times$ 在正文中已出现两到三次）。第690行把稳定性归因于“pretrained representations and synthetic exchange”，但没有相应实验；“finer label ordering”又一次把 Pearson 当成排序指标（条目25）。最后一句是附加的，与本段主题无关，且没有反映服务器 Pearson 下降等混合结果。讨论应增加的内容，例如与已有脑龄研究 MAE 的比较、结果适用的边界，在这两段中缺失。}
\conf{中}
\normal{讨论中简要回顾结果属正常；但本节几乎没有超出结果部分的新信息。}
\advice{精简；补充与外部文献的比较和适用边界。}
\revise{\authfill{与已发表脑龄研究（相近年龄范围、T1w、相近分辨率）的 MAE 比较及差异原因}；第690行末句改为：“Registration and label standardization each reduced MAE in the healthy-controls cohort, with mixed effects on the correlation-based metrics (Section~5.6).”}
\end{finding}

%---------------------------------------------------------------------
\begin{finding}{局限性部分遗漏了若干关键限制}
\loc{第692行。}
\orig{We also note several limitations ... (v) We adopt the noise strength [leiblStablePrivacy2024] to mitigate the risk that synthetic data expose some features of the original data, but real-world deployments would still require additional measures ...}
\ptype{6 信息不足；8 方法与边界说明。}
\basis{现有五条局限有实际内容。但以下影响结论可信度的限制没有列出：单次运行、无显著性检验；M4Fed 只有 2 个客户端；实验模型规模远小于“foundation-scale”，1B 参数的情形只是推算；没有隐私评估，且合成数据附带元数据文本；训练/测试划分单位不明；基线复现与原方法存在差异；跨站点评估只评估了 M4Fed。(v) 所引文献针对人脸图像（条目16）。}
\conf{中}
\normal{局限性部分不必面面俱到，但上述几项与正文的核心结论直接相关。}
\advice{补充。}
\revise{(vi) All results come from single runs without significance testing, and M4Fed was evaluated with two clients. (vii) The domain and generative models used here have fewer than 10M trainable parameters; the communication advantage for foundation-scale models is extrapolated, not measured. (viii) We did not evaluate the privacy of the exchanged synthetic volumes and their conditioning text. \authfill{其余确认后的局限}}
\end{finding}

%---------------------------------------------------------------------
\begin{finding}{“实践意义”段与结论重复引言，且表述超出证据}
\loc{第694行；第703行。}
\orig{(第694行) Exploiting all of them together is hindered by three barriers, namely ... By combining federated consolidation of siloed and public data, cloud-edge collaboration via parameter-efficient fine-tuning, and Synthetic-data-driven aggregation, M4Fed alleviates these three barriers simultaneously. Therefore, M4Fed offers a feasible route for hospitals ...\\
(第703行) ... reconciles three resources previously difficult to consolidate under realistic clinical constraints ... by combining an imputed multi-modal generative-data-mediated federated protocol with QLoRA-based cloud-edge collaboration. M4Fed can simultaneously alleviate three practical challenges ... M4Fed can simultaneously improve predictive accuracy, reduce communication overhead, and lower edge-side compute. Future work includes extending ..., balancing ..., evaluating ..., exploring hybrid aggregation ...}
\ptype{6 重复；3 放大；2 机械重复；5 修饰语堆叠。}
\basis{(a) 第694行几乎逐句复述了第102--111行（三资源、三障碍、三设计），没有新信息。(b) “alleviates these three barriers simultaneously”“offers a feasible route for hospitals”“under realistic clinical constraints”与以下事实不符：实验是单机模拟（局限 iii），默认配置下通信量增加，计算开销没有实测，也没有做隐私评估。(c) “imputed multi-modal generative-data-mediated federated protocol”连续堆叠五个修饰语；“simultaneously”在两句中重复出现；未来工作列表缺少“and”，且与局限部分重复。(d) “Synthetic-data-driven”句中误用大写。}
\conf{高}
\normal{结论简要复述贡献属正常；但第694行与第703行合计重复了第三次以上，且结论强于正文证据。}
\advice{删除或大幅精简第694行；改写第703行。}
\revise{(第694行) 建议删除；若保留，压缩为一句：“For a hospital that holds private data and has access to public data and pretrained models, M4Fed illustrates one way to combine these resources under privacy, communication, and compute constraints; validation in a real multi-institutional deployment remains future work.”\\
(第703行) We presented M4Fed, a multi-source federated learning framework that combines private data, public data, and pretrained medical models through generative-data exchange and QLoRA-based client fine-tuning. On brain-age estimation in OpenMind, the M4Fed server achieved lower MAE than FedAvg, FedProx, FAN, and GPT-FL, and QLoRA reduced client-side trainable parameters by 81.9\%; the per-round upload was smaller than that of FedAvg only for synthesis ratios below about 0.018 with the compact models used here. Future work includes additional MRI sequences and tasks, load balancing between server and clients, deployment across real institutions, and hybrid aggregation schemes.}
\end{finding}

%---------------------------------------------------------------------
\begin{finding}{全文模板化过渡与句式重复（汇总）}
\loc{第206、480行（It is worth noting that）；第455、643行（Interestingly）；第517、574行（We see two plausible reasons）；第455、517、643、645行（A possible / plausible explanation is that）；第653行（连续两个 Moreover）；第106、200行（privacy-friendly）；第694、703行（simultaneously）；几乎每段的 (i)/(ii)/(iii) 编号。}
\orig{见原文位置所列各处。}
\ptype{2 机械重复；5 表达不自然。}
\basis{这些句式本身都是正常的学术表达。问题在于它们在结果讨论中按固定模式反复出现（结果 → “two plausible reasons (i)(ii)” → 总结），使各小节读起来结构雷同，论证的推进感较弱；同时，大量编号列举把推测和观察放在同一层级上呈现。}
\conf{中}
\normal{是，单独每一处都属正常表达；只有在整体密度上才构成写作问题。}
\advice{精简；每个小节只保留一个最有依据的解释，其余删去或移入讨论。}
\revise{按第三部分“统一修改原则”处理，不逐句给出改写。}
\end{finding}

%---------------------------------------------------------------------
\begin{finding}{补充材料中的语法与表述问题}
\loc{S 第58、64、137、147、176行。}
\orig{(S58) M4Fed applies this reparameterization to both the domain model backbone and the diffusion U-Net at client nodes replacing full-model.\\
(S137) ... a fourfold spatial downsampling per axis that renders 3D diffusion tractable on hospital-grade accelerators.\\
(S147) ... which is the adaptation needed for federated personalization.\\
(S176) FAN is a decentralized baseline, architecturally aligned with augmentation-driven federated designs ...}
\ptype{5 语法与表达；9 术语一致性；1 空泛。}
\basis{S58 句尾“replacing full-model”不合语法。S137“tractable on hospital-grade accelerators”没有依据（实验使用 A100）。S147“federated personalization”在正文中从未出现。S176“architecturally aligned with augmentation-driven federated designs”含义空泛。S64 中符号 $r$ 同时表示轮次、秩和 Pearson 相关，属三重重载，建议把 Pearson 相关改记为 $r_{\mathrm{P}}$ 或在正文中写全称。}
\conf{中}
\normal{部分属轻微语言问题。}
\advice{改写。}
\revise{(S58) At client nodes, M4Fed applies this reparameterization to both the domain-model backbone and the diffusion U-Net instead of full fine-tuning.\\
(S137) ... a fourfold spatial downsampling per axis, which reduces the memory cost of 3D diffusion.\\
(S147) ... which adapts the generative model to each node's data distribution.\\
(S176) FAN is a decentralized baseline in which three clients ...}
\end{finding}

%=====================================================================
\section{综合修改}

\subsection{优先级最高、最需要修改的段落}
\begin{enumerate}[leftmargin=*]
\item \textbf{摘要（第84行）、引言末段（第141行）、表1、第694行、结论（第703行）}：为精度、通信、计算三项结论加上成立条件（服务器/客户端、$\rho_f$ 取值、只统计可训练参数），删除第694行的重复复述（条目1、2、7、40）。
\item \textbf{默认配置结果不一致（表5、表7、表8）}：先核实再改写，因为这一问题影响第5.2--5.5节所有基于单次差异的结论（条目27）。
\item \textbf{隐私表述（第202、275行）}：删除“can be exchanged freely”，核实合成数据附带元数据的内容，说明隐私未经评估（条目16）。
\item \textbf{主比较（第448行）和组件消融（第455行）}：说明比较单位不对等；核实三个消融结果几乎相同的原因；“confirming”改为相称的措辞（条目24、26）。
\item \textbf{健康状态消融（第480--482行）}：删去基于年龄标准差约 1 年的子集给出的部署建议，并调整段落顺序（条目28）。
\item \textbf{跨站点泛化（第602--604行）、噪声强度结论（第645行）、讨论第686行}：删去或限定超出证据的推论（条目32、34、37）。
\item \textbf{相关工作第159行与表1的矛盾；式(2)与第4.4节的矛盾}（条目6、14）。
\end{enumerate}

\subsection{反复出现的表达问题与统一修改原则}
\begin{enumerate}[leftmargin=*]
\item \textbf{结论与条件绑定。}凡给出百分比或“优于”的地方，同一句里写明节点（server/client）、配置（$\rho_f$、队列）和统计方式（单次运行、客户端均值）。
\item \textbf{分清三类陈述。}直接观察用过去时和具体数字；文献结论加引用并只陈述文献实际研究的对象；推测统一用“may / we hypothesize / we did not test”，且不在推测之后接“Therefore”。
\item \textbf{每个结果小节只保留一个最有依据的解释。}避免“two plausible reasons (i)(ii)”的固定模板；未经验证的解释移到讨论部分，或标注为未检验。
\item \textbf{“三元”框架只完整陈述一次。}在引言中完整陈述一次，其后以“the three resources / constraints”指代，删除第694行，结论只保留与实验结果对应的部分。
\item \textbf{术语统一。}Pearson $r$ 称为 linear correlation，不称 ranking；“healthy-controls cohort”全文统一；“granularity”“privacy-friendly”删除或给出定义；“foundation model”在指代 9.25M 模型时加以说明；方法名全称与缩写对应。
\item \textbf{拆分超长句和超长段落。}一句只承担一个论点（第106、156、200、275行）。
\item \textbf{连接词按逻辑关系使用。}“Therefore”“On the contrary”“Consequently”只在存在相应逻辑关系时使用。
\end{enumerate}

\subsection{需要补充数据、文献或与上下游核对的判断}
\begin{longtable}{p{0.04\linewidth}p{0.55\linewidth}p{0.33\linewidth}}
\toprule
\# & 需补充或核实的内容 & 关联条目 \\
\midrule
1 & 主比较、消融、敏感性分析分别使用哪个队列 & 1, 27 \\
2 & 表5、表7、表8 中默认配置结果为何不同；运行次数、随机种子；重复实验均值与标准差 & 24, 26, 27, 29 \\
3 & 三个消融变体的预测分布（是否退化为近常数预测） & 26 \\
4 & 训练/测试划分单位（受试者或扫描），以及是否存在受试者跨集合 & 20 \\
5 & 合成数据附带的条件文本是否来自源受试者元数据；隐私评估方案 & 16 \\
6 & SwinMTL 推断时是否使用元数据（式(2)） & 14 \\
7 & 各节点与数据块的对应关系；FAN、GPT-FL 复现与原方法的差异；基线通信轮数与收敛情况 & 23, 20 \\
8 & 双向、多轮总通信量；客户端显存、耗时（含 DDIM 生成） & 2, 21, 35 \\
9 & 元数据各字段缺失率；TabPFN 补全准确性；TabPFN 上下文是否包含测试行 & 18 \\
10 & 健康状态消融各测试切片的样本量与年龄分布 & 28 \\
11 & 基线在 $\mathcal{D}^{\mathrm{cross}}$ 上的结果 & 32 \\
12 & Dirichlet $\alpha=1.0$ 时 Pearson 为负是否源于训练异常 & 29 \\
13 & 表9 所用的队列与噪声强度 & 33 \\
14 & 训练细节：epoch、优化器、DDIM $\eta$/guidance、U-Net 适配器秩、$n_s^{\mathrm{priv}}$ & 20 \\
15 & VAE checkpoint 来源（正文写 BraTS/MSD，补充材料写 brain MRI generative modeling） & 30, 42 \\
16 & 脑龄估计相关工作及与已发表 MAE 的比较；3D 脑 MRI 扩散生成方法的比较 & 12, 38 \\
17 & 方法名全称与“M4Fed”的对应 & 3 \\
18 & 引用待核实（仅核对了 bib 题名，未检索全文）：leiblStablePrivacy2024（人脸图像，能否支持 MRI 去识别）；sastryDiffAug2024（$\sigma$ 范围能否用于回归）；ren2025fedfm（“now standard in medical image analysis”）；hollmann2025tabpfn（“clinical tables”）；kocak2026undertrained（放射组学样本量能否支持“单中心队列过小”）；rajabli2025brainage、leonardsen2022deep（预处理结论的具体内容）；GPT-FL 原文是否使用公开数据；OpenMind 原文中的“约 800 个子数据集”和“$113{,}921$ 条记录” & 4, 6, 11, 16, 19, 22, 30 \\
\bottomrule
\end{longtable}

\subsection{关于本审查结论性质的说明}
本报告指出的模板化过渡、句式重复、修饰语堆叠、结论放大等写作特征，在人工写作、非母语写作、多作者合稿以及使用编辑工具的写作中都很常见。\textbf{这些写作特征不能证明文本由 AI 生成}，本报告也没有、不应该据此推断文本来源或给出任何生成概率。本报告的目的仅在于指出影响论证清晰度、证据对应关系和文内一致性的具体位置，并给出修改方向。建议修改文本中凡涉及原文没有提供的信息，均以“待作者补充”标出，没有新增事实、数据或文献。

\end{document}
